% math0 manuscript. Build with: latexmk -pdf paper.tex
\documentclass[11pt]{amsart}
\usepackage[margin=1.1in]{geometry}
\usepackage{amsmath,amssymb,amsthm,mathtools}
\usepackage{url}
\usepackage{hyperref}
\usepackage{xcolor}
\usepackage{listings}
\usepackage{newunicodechar}

\lstdefinelanguage{lean}{
  keywords={theorem,lemma,def,abbrev,example,fun,by,have,show,from,let,in,if,then,else,match,with,
    namespace,end,open,import,section,variable,instance,structure,where,calc,obtain,intro,exact,
    refine,rcases,rintro,simp,omega,decide,norm_num,linarith,nlinarith,positivity,induction,cases,
    private,noncomputable,ring,push_cast},
  sensitive=true, comment=[l]{--}, morecomment=[s]{/-}{-/}, string=[b]"}
\lstset{language=lean, basicstyle=\ttfamily\footnotesize, columns=fullflexible, keepspaces=true, breaklines=true,
  keywordstyle=\color{blue!60!black}\bfseries, commentstyle=\color{gray}, frame=single,
  extendedchars=true, inputencoding=utf8,
  literate={ℕ}{{$\mathbb{N}$}}1 {ℤ}{{$\mathbb{Z}$}}1 {∀}{{$\forall$}}1 {∃}{{$\exists$}}1
    {↔}{{$\leftrightarrow$}}1 {¬}{{$\neg$}}1 {∧}{{$\wedge$}}1 {∨}{{$\vee$}}1 {→}{{$\to$}}1
    {≤}{{$\le$}}1 {⟨}{{$\langle$}}1 {⟩}{{$\rangle$}}1 {·}{{$\cdot$}}1}
% Unicode commonly found in Lean sources.
\newunicodechar{ℕ}{\ensuremath{\mathbb{N}}}
\newunicodechar{ℤ}{\ensuremath{\mathbb{Z}}}
\newunicodechar{ℚ}{\ensuremath{\mathbb{Q}}}
\newunicodechar{ℝ}{\ensuremath{\mathbb{R}}}
\newunicodechar{ℂ}{\ensuremath{\mathbb{C}}}
\newunicodechar{∀}{\ensuremath{\forall}}
\newunicodechar{∃}{\ensuremath{\exists}}
\newunicodechar{→}{\ensuremath{\to}}
\newunicodechar{←}{\ensuremath{\leftarrow}}
\newunicodechar{↔}{\ensuremath{\leftrightarrow}}
\newunicodechar{↦}{\ensuremath{\mapsto}}
\newunicodechar{¬}{\ensuremath{\neg}}
\newunicodechar{∧}{\ensuremath{\wedge}}
\newunicodechar{∨}{\ensuremath{\vee}}
\newunicodechar{≤}{\ensuremath{\le}}
\newunicodechar{≥}{\ensuremath{\ge}}
\newunicodechar{≠}{\ensuremath{\ne}}
\newunicodechar{∈}{\ensuremath{\in}}
\newunicodechar{∉}{\ensuremath{\notin}}
\newunicodechar{⊆}{\ensuremath{\subseteq}}
\newunicodechar{∩}{\ensuremath{\cap}}
\newunicodechar{∪}{\ensuremath{\cup}}
\newunicodechar{∑}{\ensuremath{\sum}}
\newunicodechar{∏}{\ensuremath{\prod}}
\newunicodechar{⟨}{\ensuremath{\langle}}
\newunicodechar{⟩}{\ensuremath{\rangle}}
\newunicodechar{·}{\ensuremath{\cdot}}
\newunicodechar{∘}{\ensuremath{\circ}}
\newunicodechar{⁻}{\ensuremath{^{-}}}
\newunicodechar{¹}{\ensuremath{^{1}}}
\newunicodechar{²}{\ensuremath{^{2}}}
\newunicodechar{₀}{\ensuremath{_{0}}}
\newunicodechar{₁}{\ensuremath{_{1}}}
\newunicodechar{₂}{\ensuremath{_{2}}}
\newunicodechar{α}{\ensuremath{\alpha}}
\newunicodechar{β}{\ensuremath{\beta}}
\newunicodechar{ε}{\ensuremath{\varepsilon}}
\newunicodechar{δ}{\ensuremath{\delta}}
\newunicodechar{π}{\ensuremath{\pi}}
\newunicodechar{σ}{\ensuremath{\sigma}}
\newunicodechar{φ}{\ensuremath{\varphi}}
\newunicodechar{∞}{\ensuremath{\infty}}
\newunicodechar{⌊}{\ensuremath{\lfloor}}
\newunicodechar{⌋}{\ensuremath{\rfloor}}
\newunicodechar{⌈}{\ensuremath{\lceil}}
\newunicodechar{⌉}{\ensuremath{\rceil}}
\newunicodechar{∣}{\ensuremath{\mid}}
\newunicodechar{≡}{\ensuremath{\equiv}}
\newunicodechar{×}{\ensuremath{\times}}
\newunicodechar{⧸}{\ensuremath{/}}
\newunicodechar{‖}{\ensuremath{\|}}
\newunicodechar{√}{\ensuremath{\surd}}
\newunicodechar{∅}{\ensuremath{\emptyset}}
\newunicodechar{⦃}{\ensuremath{\{\!|}}
\newunicodechar{⦄}{\ensuremath{|\!\}}}
\newunicodechar{𝓝}{\ensuremath{\mathcal{N}}}
\newunicodechar{≫}{\ensuremath{\gg}}
\newunicodechar{≪}{\ensuremath{\ll}}

\theoremstyle{plain}
\newtheorem{theorem}{Theorem}
\newtheorem{lemma}[theorem]{Lemma}
\newtheorem{proposition}[theorem]{Proposition}
\newtheorem{corollary}[theorem]{Corollary}
\newtheorem{conjecture}[theorem]{Conjecture}
\theoremstyle{definition}
\newtheorem{definition}[theorem]{Definition}
\newtheorem{remark}[theorem]{Remark}

\emergencystretch=3em
% Lean identifiers in running text (breakable at dots and underscores).
\expandafter\def\expandafter\UrlBreaks\expandafter{\UrlBreaks\do\_\do\.}
\DeclareUrlCommand\decl{\urlstyle{tt}}
\newcommand{\N}{\mathbb{N}}
\newcommand{\Z}{\mathbb{Z}}
\newcommand{\T}{\mathsf{T}}
\newcommand{\pr}{\operatorname{prime}}

\title[Every odd prime is a side of a Heronian triangle]{Every odd prime is a side length of a triangle with integer sides and integer area (OEIS A229159)}
\author{math0}
\address{Produced by Claude (an AI model developed by Anthropic) in the math0 project}
\date{October 18, 2026}

\begin{document}
\begin{abstract}
We prove that for every prime $p>2$ there is a non-degenerate triangle with integer side lengths and integer area one of whose sides has length $p$; this is the conjecture recorded by Michel Lagneau in 2013 in the OEIS entry A229159, whose terms are the least such areas.
The conjecture is an immediate consequence of the classical fact, attributed to Pythagoras by Proclus, that every odd $n=2m+1\ge 3$ is a leg of the right triangle with sides $n$, $(n^2-1)/2$, $(n^2+1)/2$, whose area $m(m+1)(2m+1)$ is an integer.
The result is therefore known, and we classify this note as a first formal proof of a known result, not as new mathematics.
The theorem is verified in Lean~4 against a statement that math0 formalized from the OEIS entry and that was reviewed for fidelity and frozen before the proof attempt; the definitions in that statement file are checked by the kernel against the first $19$ listed terms, and the same file shows that five later listed terms are not minimal.
\end{abstract}
\maketitle

\section{Introduction}

A \emph{Heronian triangle} is a non-degenerate triangle whose side lengths and area are integers.
The OEIS entry A229159 \cite{A229159}, submitted by Michel Lagneau on September~17, 2013, has the name ``Smallest integer areas of integer-sided triangles where at least one side is of length prime(n)'', with offset~$1$; its first terms are
\[
  0,\ 6,\ 6,\ 42,\ 66,\ 24,\ 36,\ 114,\ 966,\ 60,\ 930,\ 114,\ 126,\ 1290,\ 4230,\ 90,\ 1770,\ 330,\ 2814,\ \dots,
\]
where $a(1)=0$ records that no such triangle has a side of length $\pr(1)=2$.
The entry states the conjecture that every prime $p>2$ is a side length of a Heronian triangle (Section~\ref{sec:statement}), that is, $a(n)>0$ for all $n\ge2$.
A companion entry by the same author, A226453 \cite{A226453}, lists the integer areas of integer-sided triangles with at least one side of prime length.

\subsection*{Prior work}
The conjecture follows at once from a classical construction.
Proclus, in his commentary on Proposition~I.47 of Euclid's \emph{Elements} \cite{Proclus}, attributes to Pythagoras a rule that takes an odd number $n$ as the shorter leg of a right triangle, $(n^2-1)/2$ as the other leg and $(n^2+1)/2$ as the hypotenuse; see also Eves \cite[Sec.~3-4]{Eves}.
The lemma that precedes Proposition~X.29 of the \emph{Elements} \cite[vol.~3]{Heath} constructs, in general, two square numbers whose sum is a square; in modern terms it contains this rule as a special case.
For an odd prime $p$ the rule gives a right triangle with integer sides whose area $p(p^2-1)/4$ is an integer, since $8\mid p^2-1$.

The rule, and the fact that every odd prime is a leg of a Pythagorean triangle, were recorded in the OEIS before Lagneau's entry.
In A046079 \cite{A046079} (number of Pythagorean triangles with leg $n$), Ant King noted on January~29, 2011 that for $n>2$ the positions of the ones are the primes and their doubles, so each odd prime is a leg of exactly one Pythagorean triangle; and Richard R.~Forberg noted on May~31, 2013, with $L$ the length of the longest leg and $H$ the hypotenuse, that ``For odd $n$: $L=(n^2-1)/2$ and $H=L+1$''.
In a related direction, A042965 \cite{A042965} records (in a comment by Sharon Sela, January~25, 2002) that its terms are the nonnegative integers that are differences of two squares, and N.~J.~A.~Sloane added on May~30, 2014, after Lagneau's entry, that the primes in the sequence are all the odd primes.

The A229159 entry (revision \#18 of September~16, 2017, which is the version we read) still states the conjecture and does not mention this construction.
The math0 referees found no source that addresses the conjecture of A229159 directly.
Mathlib \cite{Mathlib} contains the classification of Pythagorean triples (\path{Mathlib/NumberTheory/PythagoreanTriples.lean}) but not this statement, and the statement is not among those of the formal-conjectures repository \cite{FC}.
The mathematics in this note is classical and elementary.
What we add is a complete proof that matches a kernel-checked Lean proof of a frozen formal statement (Section~\ref{sec:lean}); to our knowledge it is the first formal proof of the statement.
We also record (Remark~\ref{rem:errata}) a by-product of the check of the formal definitions against the OEIS data: the listed terms $a(20)$, $a(25)$, $a(29)$, $a(32)$ and $a(40)$ are not minimal.

\section{Statement}\label{sec:statement}

\subsection*{The OEIS conjecture}
The comments of A229159 contain the following sentence, verbatim:
\begin{quote}
Conjecture: for all prime p $>$ 2 there exists an integer-sided triangle with integer area where at least one side is of length p.
\end{quote}
This is the wording of the current revision (\#18).
Michel Lagneau added the conjecture when he submitted the entry on September~17, 2013 (revision \#7 in the entry's edit history), in the form ``Conjecture: for all prime p $>$ 2 there exists an integer-sided triangles where at least one side is of p length.''
The words ``with integer area'' and the present phrasing come from a copy-edit by Jon E.~Schoenfield of September~15, 2017 (revision \#12, approved as revision \#18).
Read together with the entry's name, which asks for integer areas, the original wording expresses the same conjecture; we formalize the current wording.

\subsection*{Mathematical formulation}
By Heron's formula, a triangle with side lengths $x,y,z$ has area $S$ with
\[
  16S^2=(x+y+z)(-x+y+z)(x-y+z)(x+y-z).
\]
A triple of positive reals is the triple of side lengths of a non-degenerate triangle if and only if it satisfies the three strict triangle inequalities, and the right-hand side is then positive.
This motivates the following definition.

\begin{definition}\label{def:T}
For $x,y,z,A\in\N$ write $\T(x,y,z;A)$ if $x<y+z$, $y<x+z$, $z<x+y$ and
\[
  16A^2=(x+y+z)(-x+y+z)(x-y+z)(x+y-z)\qquad\text{in }\Z.
\]
\end{definition}

Thus $\T(x,y,z;A)$ holds exactly when $x,y,z$ are the side lengths of a non-degenerate triangle whose area is the natural number $A$; in that case $x,y,z\ge1$ and $A\ge1$.

\begin{theorem}\label{thm:main}
For every prime $p>2$ there are $x,y,z,A\in\N$ with $\T(x,y,z;A)$ and $p\in\{x,y,z\}$.
\end{theorem}

The proof (Section~\ref{sec:proof}) gives an explicit witness: if $p=2m+1$, then
\[
  \T\bigl(2m+1,\ 2m^2+2m,\ 2m^2+2m+1;\ m(m+1)(2m+1)\bigr).
\]

\subsection*{Lean statement}
The statement is the definition \decl{Math0.Statements.A229159.statement} in the frozen math0 statement file \path{lean/Math0/Statements/OEIS/A229159.lean}.
The two definitions it consists of are, verbatim (docstrings omitted; both lie in the namespace \decl{Math0.Statements.A229159}):
\begin{lstlisting}
def IsIntTriangle (x y z A : ℕ) : Prop :=
  x < y + z ∧ y < x + z ∧ z < x + y ∧
    16 * (A : ℤ) ^ 2 =
      ((x : ℤ) + y + z) * (-(x : ℤ) + y + z) * ((x : ℤ) - y + z) * ((x : ℤ) + y - z)

def statement : Prop :=
  ∀ p : ℕ, p.Prime → 2 < p →
    ∃ x y z A : ℕ, IsIntTriangle x y z A ∧ (x = p ∨ y = p ∨ z = p)
\end{lstlisting}
\texttt{IsIntTriangle x y z A} is $\T(x,y,z;A)$, and \decl{statement} is Theorem~\ref{thm:main}.

\subsection*{Encoding}
An ``integer-sided triangle'' is a triple of natural numbers satisfying the strict triangle inequalities, so degenerate triangles are excluded and every side is at least~$1$.
Without this exclusion the statement would be trivial, since $(p,p,2p)$ would be a ``triangle'' of area~$0$, and $a(1)=0$ would no longer mean that there is no triangle.
``Integer area'' means that the area given by Heron's formula is a natural number $A$; the formal statement uses Heron's formula as the definition of the area and does not refer to Euclidean geometry.
The four factors are computed in $\Z$, so truncated subtraction in $\N$ plays no role.
``At least one side is of length $p$'' is the disjunction $x=p\lor y=p\lor z=p$.
The entry's Maple program uses the same notion: for sides $a\le b\le c$ it computes $x=s(s-a)(s-b)(s-c)$ with $s=(a+b+c)/2$, takes $\sqrt{x}$ only when $x>0$, and tests whether it is an integer; since $16x$ is the Heron product, this is the condition $\T(a,b,c;\sqrt{x})$.

\subsection*{Check against the OEIS data}
The statement file also defines the sequence itself:
\begin{lstlisting}
def areas (m : ℕ) : Set ℕ :=
  {A | ∃ x y z : ℕ, IsIntTriangle x y z A ∧ (x = m ∨ y = m ∨ z = m)}

noncomputable def prime (n : ℕ) : ℕ := Nat.nth Nat.Prime (n - 1)

noncomputable def a (n : ℕ) : ℕ := sInf (areas (prime n))
\end{lstlisting}
Here \texttt{prime n} is the $n$th prime with \texttt{prime 1 = 2}, and \decl{sInf} on $\N$ is the least element of a nonempty set and $0$ for the empty set, which matches the convention $a(1)=0$.
In these terms \decl{statement} says that \texttt{areas p} is nonempty for every prime $p>2$, that is, $a(n)\ne0$ for $n\ge2$.

The file proves \decl{a_data}, which states that $(a(1),\dots,a(19))$ equals the first $19$ listed terms (the primes up to~$67$).
The value $a(1)=0$ is proved by hand from \texttt{areas 2 =~}$\emptyset$ (Remark~\ref{rem:two}).
For $n=2,\dots,19$ a single \texttt{decide +kernel} evaluates, in the kernel, a Boolean certificate for each table entry $(n,p,A_0,b,c)$: that $p$ is prime and is the $n$th prime (via \decl{Nat.count}); that $(p,b,c)$ satisfies the triangle inequalities with Heron product $16A_0^2$ (the witnesses are the triangles in the entry's table); and that no triangle with a side $p$ has an integer area below $A_0$.
For the last check, a triangle $(p,b,c)$ with $b\le c$ has $d=c-b<p$, and for fixed $p$ and $d$ the Heron product of $(p,b,b+d)$ is monotone in $b$ (lemma \decl{heronN_mono}); so for each $d<p$ the check scans $b=0,1,2,\dots$ until the Heron product reaches $16A_0^2$ and verifies that no earlier $(p,b,b+d)$ satisfying the triangle inequalities has a Heron product that is $16$ times a square, certifying each non-square $N$ by an $r$ with $r^2<N<(r+1)^2$.
Soundness of these helpers is proved in the file, so \decl{a_data} is a statement about the noncomputable \decl{a} above.
Independently of the formal check, a brute-force search by the math0 fidelity referee agrees with all $48$ listed terms except the five discussed next.

\begin{remark}[Listed terms that are not minimal]\label{rem:errata}
The statement file proves, by kernel evaluation of witness certificates,
\begin{lstlisting}
theorem oeis_errata : a 20 ≤ 11928 ∧ a 25 ≤ 396 ∧ a 29 ≤ 330 ∧ a 32 ≤ 6288 ∧ a 40 ≤ 2970
\end{lstlisting}
using the triangles $(71,505,560)$, $(97,11,90)$, $(109,11,100)$, $(131,265,390)$ and $(173,36,187)$, whose integer areas are $11928$, $396$, $330$, $6288$ and $2970$.
The listed terms are $a(20)=14910$, $a(25)=420$, $a(29)=840$, $a(32)=11790$ and $a(40)=3276$, so they are not minimal.
We do not determine the true values here, and we note only that the entry's Maple program limits the sides to $500$ and stops at the first triangle it finds, which is in general not one of least area.
These errata concern the listed values only, not the conjecture, which is an existence statement.
\end{remark}

\section{Proof}\label{sec:proof}

The proof has one lemma, which is the Lean lemma \decl{isIntTriangle_odd}.

\begin{lemma}\label{lem:odd}
Let $m\ge1$ be an integer and put
\[
  x=2m+1,\qquad y=2m^2+2m,\qquad z=2m^2+2m+1,\qquad A=m(m+1)(2m+1).
\]
Then $\T(x,y,z;A)$.
\end{lemma}

\begin{proof}
\emph{Triangle inequalities.}
We have $y+z-x=4m^2+2m>0$, so $x<y+z$.
Next $y<z<x+z$.
Finally $x+y-z=2m>0$, so $z<x+y$.
The first and last inequalities use $m\ge1$.

\emph{Heron's identity.}
Computing in $\Z$,
\begin{align*}
  x+y+z &= 4m^2+6m+2 = 2(m+1)(2m+1), & -x+y+z &= 4m^2+2m = 2m(2m+1),\\
  x-y+z &= 2m+2 = 2(m+1), & x+y-z &= 2m.
\end{align*}
The product of the four factors is
\[
  2(m+1)(2m+1)\cdot 2m(2m+1)\cdot 2(m+1)\cdot 2m = 16\,m^2(m+1)^2(2m+1)^2 = 16A^2. \qedhere
\]
\end{proof}

In Lean, the inequality $x<y+z$ is closed by \texttt{nlinarith}, the other two by \texttt{omega} (which treats $m^2$ as an atom), and Heron's identity by \texttt{push\_cast} followed by \texttt{ring}, which proves it as a polynomial identity in $m$ over $\Z$.

\begin{proof}[Proof of Theorem~\ref{thm:main}]
Let $p$ be a prime with $p>2$.
Then $p\ne2$, so $p$ is odd and $p=2m+1$ for some $m\in\N$; since $p>2$ we have $m\ge1$.
Lemma~\ref{lem:odd} gives $\T(p,\,2m^2+2m,\,2m^2+2m+1;\,m(m+1)(2m+1))$, and the first side is $p$.
\end{proof}

In Lean, the oddness of $p$ is \decl{Nat.Prime.odd_of_ne_two}, and the witness is supplied with the first disjunct $x=p$.

\begin{remark}
The triangle of Lemma~\ref{lem:odd} is the right triangle of the Pythagorean rule: $x^2+y^2=z^2$, and $A=xy/2$.
Neither fact is used in the proof, which checks Heron's identity directly.
Primality is used only to obtain that $p$ is odd, so the argument shows that every odd $n\ge3$ is a side length of a Heronian triangle.
It also gives $1\le a(n)\le \pr(n)\bigl(\pr(n)^2-1\bigr)/4$ for $n\ge2$; we did not state this bound in Lean.
The witness is in general far from minimal: for $p=7$ it has area $84$, whereas $a(4)=42$, attained by $(7,15,20)$.
\end{remark}

\begin{remark}\label{rem:two}
The hypothesis $p>2$ is necessary; the statement file proves \texttt{areas 2 =~}$\emptyset$.
By symmetry of $\T$ it suffices to consider triangles $(2,b,c)$ with $b\le c$, and then $c<b+2$.
If $c=b$, then $2<2b$ gives $b\ge2$, and the Heron product is $(2b+2)(2b-2)\cdot2\cdot2=16(b^2-1)$; but $b^2-1$ is not a square, since $(b-1)^2<b^2-1<b^2$.
If $c=b+1$, the Heron product is $3(2b+3)(2b-1)$, which is odd and so is not $16A^2$.
\end{remark}

\section{Machine verification}\label{sec:lean}

\subsection*{Files and declarations}
The proof is in the file \path{lean/Math0/Frontier/OEIS/A229159.lean} of the math0 repository, \url{https://github.com/Foundation-0/math0}.
It imports the frozen statement file \path{lean/Math0/Statements/OEIS/A229159.lean} (Section~\ref{sec:statement}; SHA-256 \texttt{634d66ee\ldots481a}, unchanged since the fidelity review) and the audit commands in \path{lean/Math0/Audit.lean}.
The main declaration is \decl{Math0.Frontier.A229159.result}, of type \decl{Math0.Statements.A229159.statement}; its only auxiliary lemma is the private theorem \decl{isIntTriangle_odd}.
The proof file is, verbatim except for the license header, the imports and the docstrings:
\begin{lstlisting}
namespace Math0.Frontier.A229159
open Math0.Statements.A229159

private theorem isIntTriangle_odd (m : ℕ) (hm : 1 ≤ m) :
    IsIntTriangle (2 * m + 1) (2 * m ^ 2 + 2 * m) (2 * m ^ 2 + 2 * m + 1)
      (m * (m + 1) * (2 * m + 1)) := by
  refine ⟨?_, by omega, ?_, ?_⟩
  · nlinarith
  · omega
  · push_cast
    ring

theorem result : statement := by
  intro p hp h2
  obtain ⟨m, rfl⟩ : Odd p := hp.odd_of_ne_two (by omega)
  exact ⟨_, _, _, _, isIntTriangle_odd m (by omega), Or.inl rfl⟩

end Math0.Frontier.A229159

#math0_settles Math0.Frontier.A229159.result Math0.Statements.A229159.statement
\end{lstlisting}
In Mathlib, \texttt{Odd p} unfolds to $\exists m,\ p=2m+1$, so the pattern \decl{rfl} replaces $p$ by $2m+1$.
The proof file uses no \texttt{decide}, \texttt{native\_decide}, \texttt{sorry} or \texttt{set\_option}.
The kernel evaluations described in Section~\ref{sec:statement} (\decl{a_data}, \decl{oeis_errata}) live in the statement file, use \texttt{decide +kernel} (not \texttt{native\_decide}) and depend only on \decl{propext}, \decl{Classical.choice} and \decl{Quot.sound}; \decl{result} does not depend on them.

\subsection*{Audit}
The command \texttt{\#math0\_settles}, defined in \path{lean/Math0/Audit.lean} for statements that math0 formalizes itself, fails unless (1) \decl{statement} is a \texttt{Prop}-valued definition without universe parameters; (2) \decl{statement} and \decl{result} use no axioms besides \decl{propext}, \decl{Classical.choice} and \decl{Quot.sound} (this rejects \decl{sorryAx}, the axiom \decl{Lean.ofReduceBool} introduced by \texttt{native\_decide}, and custom axioms); and (3) the type of \decl{result} is syntactically the statement or its negation, up to metadata annotations, with no unfolding.
Running \texttt{lake env lean Math0/Frontier/OEIS/A229159.lean} in the directory \path{lean} prints
\begin{lstlisting}[basicstyle=\ttfamily\scriptsize]
math0 audit: 'Math0.Statements.A229159.statement' depends on axioms [propext]
math0 audit: 'Math0.Frontier.A229159.result' depends on axioms [Classical.choice, Quot.sound, propext]
math0 statement check: 'Math0.Frontier.A229159.result' proves 'Math0.Statements.A229159.statement'
\end{lstlisting}
so \decl{result} proves the statement (it does not refute it) and depends only on the three standard axioms.

\subsection*{Toolchain and reproduction}
We use Lean~4 \cite{Lean4} \texttt{v4.33.1} and Mathlib \cite{Mathlib} \texttt{v4.33.1}, pinned in \path{lean/}; the repository also pins formal-conjectures \cite{FC} at commit \texttt{c77fdd4}, which this result does not use.
To reproduce, run \path{scripts/verify.sh} from the repository root: it fetches the Mathlib build cache, builds all math0 results with \texttt{lake build} (re-running every audit command), checks that the pinned \texttt{formal\_conjectures} and \texttt{mathlib} checkouts are unmodified, and reports the number of passed statement checks.
To check this result alone, run \texttt{lake build Math0.Statements.OEIS.A229159} (about one minute, almost all of it the kernel evaluations of the data check) and then \texttt{lake env lean Math0/Frontier/OEIS/A229159.lean} in the directory \path{lean}; the second command only elaborates the short proof file, and most of its running time is spent loading Mathlib.

{\small
\begin{thebibliography}{99}

\bibitem{Eves}
H.~Eves, \emph{An Introduction to the History of Mathematics, with Cultural Connections by Jamie H. Eves}, 6th ed., Saunders College Publishing, Fort Worth, 1990.

\bibitem{FC}
M.~Firsching, P.~Lezeau, S.~Mercuri, M.~Z. Horv\'ath, Y.~Dillies, C.~S\"onne, E.~Wieser, F.~Zhang, T.~Hubert, B.~Ag\"uera y Arcas, and P.~Kohli, \emph{Formal Conjectures: An open and evolving benchmark for verified discovery in mathematics}, arXiv:2605.13171, 2026. Repository: \url{https://github.com/google-deepmind/formal-conjectures}.

\bibitem{Heath}
T.~L.~Heath, \emph{The Thirteen Books of Euclid's Elements}, translated with introduction and commentary, 2nd ed., 3 vols., Cambridge University Press, 1926; reprinted by Dover, New York, 1956.

\bibitem{Mathlib}
The mathlib Community, \emph{The Lean mathematical library}, in: Proceedings of the 9th ACM SIGPLAN International Conference on Certified Programs and Proofs (CPP 2020), ACM, 2020, 367--381.

\bibitem{Lean4}
L.~de Moura and S.~Ullrich, \emph{The Lean 4 theorem prover and programming language}, in: Automated Deduction -- CADE 28, Lecture Notes in Comput. Sci. \textbf{12699}, Springer, 2021, 625--635.

\bibitem{Proclus}
Proclus, \emph{A Commentary on the First Book of Euclid's Elements}, translated with introduction and notes by G.~R.~Morrow, Princeton University Press, Princeton, 1970; 2nd ed.\ with a foreword by I.~Mueller, 1992.

\bibitem{A042965}
The OEIS Foundation Inc., \url{https://oeis.org/A042965} (nonnegative integers not congruent to $2$ mod $4$), entry by N.~J.~A. Sloane.

\bibitem{A046079}
The OEIS Foundation Inc., \url{https://oeis.org/A046079} (number of Pythagorean triangles with leg $n$), entry by E.~W. Weisstein.

\bibitem{A226453}
The OEIS Foundation Inc., \url{https://oeis.org/A226453}, entry by M.~Lagneau, Sep.~16, 2013.

\bibitem{A229159}
The OEIS Foundation Inc., \url{https://oeis.org/A229159}, entry by M.~Lagneau, Sep.~17, 2013; revision \#18, Sep.~16, 2017; edit history at \url{https://oeis.org/history?seq=A229159}.

\end{thebibliography}
}

\end{document}
